% !TEX program = xelatex
\documentclass[10pt,a4paper]{article}

% ---- Пакеты для XeLaTeX и кириллицы ----
\usepackage{fontspec}
\usepackage{polyglossia}
\setmainlanguage{russian}
\setotherlanguage{english}

% ---- Шрифт с кириллицей ----
\setmainfont{Liberation Serif}
\setsansfont{Liberation Sans}
\setmonofont{Liberation Mono}

% ---- Поля и геометрия ----
\usepackage[top=1.5cm, bottom=1.5cm, left=1.5cm, right=1.5cm]{geometry}
\usepackage{parskip}
\setlength{\parindent}{0pt}

% ---- Продвинутые ссылки ----
\usepackage{hyperref}
\hypersetup{
    colorlinks=true,
    linkcolor=blue,
    urlcolor=blue,
    citecolor=blue
}

% ---- Полезные пакеты ----
\usepackage{enumitem}
\setlist{nosep, left=0pt}
\usepackage{tabularx}
\usepackage{array}
\usepackage{booktabs}
\usepackage{xcolor}
\usepackage{titlesec}

% ---- Настройка заголовков разделов ----
\titleformat{\section}{\Large\bfseries\uppercase}{}{0em}{}[\vspace{-0.3em}\rule{\textwidth}{0.5pt}]
\titlespacing*{\section}{0pt}{0.8em}{0.3em}

\titleformat{\subsection}{\large\bfseries}{}{0em}{}
\titlespacing*{\subsection}{0pt}{0.4em}{0.2em}

% ---- Пользовательские команды ----
\newcommand{\projectEntry}[3]{\textbf{#1} \hfill \textit{#2}\par #3\par\vspace{0.3em}}

% ---- Общие данные (имя, телефон, email и т.д.) ----
% ============================================================
% ОБЩИЕ ПЕРСОНАЛЬНЫЕ ДАННЫЕ
% ============================================================
% Меняешь здесь — меняется во всех резюме

\newcommand{\PersonName}{Кузык Иван Романович}
\newcommand{\PersonPhone}{+7 (916) 574-52-99}
\newcommand{\PersonEmail}{vanya.kuzyk1@gmail.com}
\newcommand{\PersonGitHub}{github.com/vanya890}
\newcommand{\PersonCity}{г.~Москва, м.~Бульвар Дмитрия Донского}
\newcommand{\PersonCitizenship}{РФ}
\newcommand{\PersonDrivingLicense}{Категория B}

% ============================================================
% КОМАНДА ДЛЯ ЗАГОЛОВКА РЕЗЮМЕ
% ============================================================
% Использование: \makeHeader{Должность}
\newcommand{\makeHeader}[1]{%
\begin{center}
    {\Huge\bfseries \PersonName}\\[0.3em]
    {\large #1}\\[0.5em]
    \begin{tabular}{ll}
        Телефон: & \PersonPhone \\
        Email:   & \href{mailto:\PersonEmail}{\PersonEmail} \\
        GitHub:  & \href{https://\PersonGitHub}{\PersonGitHub} \\
        Город:   & \PersonCity \\
        Гражданство: & \PersonCitizenship
    \end{tabular}
\end{center}
}

% ============================================================
% ОБЩЕЕ ОБРАЗОВАНИЕ
% ============================================================
% Использование: \makeEducation
\newcommand{\makeEducation}{%
\section{Образование}

\projectEntry{Московский политехнический университет --- Факультет информационных технологий}{2022--2026}
{Бакалавриат, направление «Информатика и вычислительная техника», специализация «Киберфизические системы». ВКР «Разработка образовательного 3D-симулятора робототехники» --- отлично. Средний балл --- 5.0.}

\projectEntry{ГБПОУ «Образовательный комплекс «Юго-Запад»}{2018--2022}
{Среднее специальное образование, специальность «Информационные системы и программирование». Диплом с отличием. Демонстрационный экзамен по компетенции «Разработка компьютерных игр и мультимедийных приложений» --- отлично.}
}

% ============================================================
% ОБЩАЯ ДОПОЛНИТЕЛЬНАЯ ИНФОРМАЦИЯ
% ============================================================
% Использование: \makeFooter{Ключевые компетенции}{Личные качества}
\newcommand{\makeFooter}[2]{%
\section{Дополнительная информация}

\begin{itemize}
    \item \textbf{GitHub:} \href{https://\PersonGitHub}{\PersonGitHub} --- #1
    \item \textbf{Ключевые компетенции:} #2
    \item \textbf{Личные качества:} Обучаемость, ответственность, умение работать с большим объёмом информации, системное мышление, ориентация на результат
    \item \textbf{Водительские права:} \PersonDrivingLicense
\end{itemize}
}

\begin{document}

% ============================================================
% ЗАГОЛОВОК
% ============================================================
\makeHeader{Инженер киберфизических систем / Специалист по робототехнике}

% ============================================================
% О СЕБЕ
% ============================================================
\section{О себе}

Инженер-разработчик в области киберфизических систем, робототехники и систем управления. Выпускник Московского политехнического университета по специализации «Киберфизические системы» (факультет информационных технологий). Имею практический опыт работы с \textbf{ROS (Robot Operating System)}, микропроцессорными системами управления, системами технического зрения, 3D-симуляцией и программированием робототехнических платформ (Clover, TurtleBot3). Разработал образовательный 3D-симулятор робототехники в качестве ВКР. Умею проектировать алгоритмы управления, обрабатывать данные с сенсоров и визуализировать результаты.

% ============================================================
% ОПЫТ РАБОТЫ (ПРОЕКТНАЯ ДЕЯТЕЛЬНОСТЬ)
% ============================================================
\section{Опыт работы}

\subsection{Проектная деятельность / Фриланс}

\projectEntry{Разработчик образовательного 3D-симулятора робототехники (ВКР)}{2025--2026}
{Спроектировал и реализовал \textbf{EduRoboSim} --- 3D-симулятор для обучения основам робототехники. Стек: \textbf{Tauri (Rust + WebView2) + React + Three.js}. Разработал физический движок с 2D OBB-коллизиями и скольжением вдоль стен. Реализовал дифференциальный привод и кинематику на плоскости XZ. Интегрировал систему сенсоров: 8 ультразвуковых дальномеров (raycast до 4 м), энкодер (интеграция линейной скорости), гироскоп. Создал собственный язык программирования EduScript (лексер, парсер AST, VM). Разработал редактор уровней с расстановкой препятствий, рамп и триггеров.}

\projectEntry{Лабораторные работы с ROS (Clover \& TurtleBot3)}{2024--2025}
{Выполнил цикл лабораторных работ в среде \textbf{ROS (Robot Operating System)}:
\begin{itemize}
    \item Разработал полётное задание для дрона \textbf{Clover} --- полёт по маркерам ArUco, фотографирование повреждений, автоматическая посадка.
    \item Реализовал алгоритм поиска отсутствующих маркеров («полыней») с оптимизацией маршрута сканирования.
    \item Разработал систему поиска маркеров в лесистой местности с использованием \textbf{лидара} для обхода препятствий.
    \item Запускал и настраивал симуляцию \textbf{TurtleBot3} в Gazebo, работал с RViz для визуализации.
\end{itemize}}

\projectEntry{Система визуальной локализации робота (Курсовая, 4 семестр)}{2024}
{Разработал desktop-приложение на \textbf{C\# Windows Forms} с использованием \textbf{OpenCvSharp} для определения положения и траектории робота по стереопотоку с двух камер. Реализовал: захват кадров, конвертацию в оттенки серого, вычисление абсолютной разницы кадров, расчёт центра масс движения, обновление позиции и угла поворота по дифференциальной модели. Визуализация траектории на карте.}

\projectEntry{Курсовые проекты по системам управления}{2023--2026}
{Выполнил серию курсовых проектов на отлично:
\begin{itemize}
    \item «Проектирование алгоритмов систем управления»
    \item «Системы технического зрения в автоматизированных системах управления»
    \item «Задачи планирования движения и навигации в робототехнике»
    \item «Технологии визуализации данных систем управления»
    \item «Программная обработка данных в системах мониторинга»
\end{itemize}}

\projectEntry{Децентрализованный P2P VPN на базе DHT}{2024--2025}
{Разработал на \textbf{C\#} децентрализованный P2P VPN для Windows. Реализовал DHT (Kademlia) для обнаружения пиров, PEX (Peer Exchange) для обмена списками узлов, шифрование трафика по протоколу \textbf{Noise (ChaCha20-Poly1305)}. Интегрировал виртуальный сетевой адаптер Wintun для создания L3-сети.}

\subsection{Практика}

\projectEntry{Техник-программист (ГБПОУ «Юго-Запад»)}{Апрель--Июнь 2022}
{Проверка работоспособности ПК, обновление и установка ОС и ПО, настройка сетевых подключений, установка принтеров и периферийного оборудования.}

\projectEntry{Оператор БД (Издательство «Просвещение»)}{Май--Июль 2021}
{Заполнение базы данных информацией, описание изображений для поиска по ключевым словам.}

% ============================================================
% ОБРАЗОВАНИЕ
% ============================================================
\makeEducation

% ============================================================
% КЛЮЧЕВЫЕ ДИСЦИПЛИНЫ
% ============================================================
\section{Ключевые дисциплины}

\begin{tabular}{p{6cm} p{10cm}}
    \textbf{Робототехника и управление:} & Проектирование алгоритмов систем управления, методы и алгоритмы управления динамическими системами, разработка систем управления БПЛА, задачи планирования движения и навигации \\
    \textbf{Компьютерное зрение:} & Системы технического зрения, нейронные сети глубокого обучения в обработке изображений, разработка систем виртуальной и дополненной реальности \\
    \textbf{Микропроцессорные системы:} & Микропроцессорные системы управления, интерфейсы и протоколы обмена данными, программное управление электроприводом \\
    \textbf{Программирование:} & Программирование и алгоритмизация, численные методы, дискретная математика \\
    \textbf{Моделирование:} & Моделирование электронных схем, моделирование систем сервисной и мобильной робототехники
\end{tabular}

% ============================================================
% НАВЫКИ
% ============================================================
\section{Навыки}

\begin{tabular}{p{3.5cm} p{12cm}}
    \textbf{Робототехника:} & ROS (Robot Operating System), Gazebo, RViz, Clover, TurtleBot3, планирование движения, навигация \\
    \textbf{Языки:} & C\#, Python, TypeScript, JavaScript, C++ (базовый) \\
    \textbf{Комп.зрение:} & OpenCV (OpenCvSharp), системы технического зрения, нейронные сети (CNN), обработка изображений \\
    \textbf{Микропроцессоры:} & Микропроцессорные системы управления, схемотехника, электротехника \\
    \textbf{3D-симуляция:} & Three.js / React Three Fiber, Unity, A-Frame \\
    \textbf{Инструменты:} & Git, Docker, Linux \\
    \textbf{Прочее:} & C\#/.NET (Windows Forms), Tauri, Matplotlib, NumPy \\
    \textbf{Языки:} & Русский --- родной, Английский --- A2 (элементарный)
\end{tabular}

% ============================================================
% ДОПОЛНИТЕЛЬНАЯ ИНФОРМАЦИЯ
% ============================================================
\makeFooter{открытые проекты по робототехнике, компьютерному зрению, симуляции}{Программирование робототехнических систем, компьютерное зрение, 3D-симуляция, системы управления, обработка сенсорных данных}

\end{document}